\begin{table}[H]
\centering
\scriptsize
\caption{Guia de decisão: qual tecnologia escolher para cada caso?}
\label{quadro:guia_decisao}
\begin{tabular}{|p{3.3cm}|p{2.3cm}|p{2.2cm}|p{2.4cm}|p{4.2cm}|}
\hline
\textbf{Caso de uso} & \textbf{Escolha recomendada} & \textbf{Alternativa} & \textbf{Evitar quando} & \textbf{Evidência experimental (este estudo)} \\ \hline
Comunicação interna entre microsserviços, alto volume e baixa latência & gRPC & REST & SOAP & Cenário ideal: 0,73 ms de latência média e 1.367 req/s (maior vazão sequencial) \\ \hline
Interação bidirecional em tempo real (chat, jogos, painéis operacionais) & WebSocket & SSE (se for só servidor→cliente) & SOAP, Webhook & Cenário ideal: 0,55 ms de latência média (menor entre as 7) \\ \hline
Notificações e feeds unidirecionais servidor→cliente (cotações, status, alertas) & SSE & WebSocket & Webhook & Cenário ideal: 6,70 ms / 147 req/s \\ \hline
API pública / integração geral entre sistemas heterogêneos & REST & GraphQL & — & Cenário ideal: 6,03 ms / 164 req/s \\ \hline
Front-end/mobile com consultas flexíveis (evitar over/underfetching) & GraphQL & REST & — & Cenário ideal: 7,93 ms (1,91 ms acima do REST) \\ \hline
Sistemas corporativos legados com contrato formal (WSDL) e WS-* & SOAP & REST & — & Cenário ideal: 5,96 ms / 166 req/s \\ \hline
Notificação de eventos a sistemas de terceiros (pagamentos, SaaS) & Webhook & SSE / REST & — & Cenário ideal: 15,83 ms / 63 req/s \\ \hline
Rede instável / móvel (alta latência e perda de pacotes) & WebSocket (menor latência) — gRPC (maior taxa de sucesso) & — & Tecnologias que abrem nova conexão a cada requisição, quando a perda é alta & Cenário crítico: WebSocket 506,52 ms; gRPC 100,00\% de sucesso \\ \hline
Infraestrutura com CPU/RAM restritas (custo de nuvem) & gRPC (menor CPU/req.) e gRPC (menor RAM) & — & Webhook (maior CPU/req.: 10,07 ms) & CPU: gRPC 0,596 ms/req. × Webhook 10,069 ms/req.; RAM pico: gRPC 73 MiB \\ \hline
\end{tabular}
{\footnotesize\\ \textbf{Fonte: Elaborada pelos autores a partir dos resultados experimentais (2026).}}
\end{table}
